\section{Backlash and Electoral Integrity}

The trust results above suggest that BVR was not generally experienced as a threatening reform, but a distinct concern remains: electoral modernization can itself become a focal point for partisan suspicion. In Brazil this concern became especially salient during the Bolsonaro presidency, when attacks on the TSE and on electronic voting helped turn election administration into an object of political conflict. If those messages stuck, respondents aligned with Bolsonaro should be less willing to say that votes are counted correctly, and that skepticism could be sharper in municipalities where biometric registration made the digital electoral infrastructure more visible.

\textbf{Data and outcomes.} The reproducible public LAPOP file available for this exercise is the Brazil 2021 wave. That file includes an electoral-integrity battery but does not include individual presidential vote recall, so I use approval of the executive as the respondent-level proxy for Bolsonaro alignment. The first outcome is a binary indicator equal to one if the respondent says votes are counted correctly \emph{always} and zero if the answer is \emph{sometimes} or \emph{never}. The second outcome is a binary indicator equal to one if the respondent says politicians can identify how each person voted \emph{never}, and zero otherwise; higher values therefore correspond to stronger belief in ballot secrecy. Because these integrity items were asked in a narrower 2021 survey module, the usable control set is more limited than in the main trust regressions. I therefore retain low education, female, white, age, urban residence, and log municipality population.

\textbf{Empirical strategy.} I estimate the following cross-sectional linear probability model:

\begin{equation}
\label{eq:lapop-bvr-backlash}
y_{i,m} = \beta_1 \text{BVR}_m + \beta_2 \text{ApproveBolsonaro}_i + \beta_3 \left(\text{BVR}_m \times \text{ApproveBolsonaro}_i\right) + \mathbf{X}'_{i,m}\Gamma + \varepsilon_{i,m},
\end{equation}

where $y_{i,m}$ is one of the two integrity outcomes for respondent $i$ in municipality $m$, $\text{BVR}_m$ equals one if the municipality had entered the biometric regime by 2021, and $\text{ApproveBolsonaro}_i$ equals one when the respondent rates the executive's performance as good or very good. The coefficient on $\text{ApproveBolsonaro}_i$ therefore captures the baseline integrity gap between Bolsonaro-aligned and other respondents in municipalities without BVR exposure, while $\beta_3$ measures whether that gap differs in municipalities already under BVR. A negative interaction would be consistent with backlash: Bolsonaro-aligned respondents in BVR municipalities would be especially skeptical. A positive interaction would imply the opposite. All specifications use LAPOP survey weights, and standard errors are clustered at the municipality level.

\begin{table}[ht]
    \caption{Backlash and electoral-integrity regressions, LAPOP Brazil 2021}
    \label{tbl:lapop-backlash-2021}
    \begin{tabular*}{\textwidth}{@{\extracolsep{\fill}}lcccc}
\doubletoprule
 & \multicolumn{2}{c}{Count Fair} & \multicolumn{2}{c}{Ballot Secret} \\
\cmidrule(lr){2-3}\cmidrule(lr){4-5}
 & (1) & (2) & (3) & (4) \\
\midrule
BVR & 0.016\\ (0.051) & -0.030\\ (0.060) & 0.028\\ (0.046) & 0.057\\ (0.051) \\
Bolsonaro approval & -0.224***\\ (0.054) & -0.324***\\ (0.058) & -0.049\\ (0.057) & 0.006\\ (0.068) \\
BVR x Bolsonaro approval &  & 0.198*\\ (0.107) &  & -0.112\\ (0.105) \\
\midrule
Interaction included & No & Yes & No & Yes \\
Observations & 698 & 698 & 715 & 715 \\
Municipalities & 338 & 338 & 376 & 376 \\
Weighted mean of dependent variable & 0.420 & 0.420 & 0.329 & 0.329 \\
$R^2$ & 0.063 & 0.070 & 0.033 & 0.035 \\
\doublebottomrule
\end{tabular*}

    \begin{minipage}{\textwidth}
        {\footnotesize \textbf{Note:} Weighted linear-probability models using the public Brazil 2021 LAPOP wave. Columns (1)--(2) use an indicator for saying votes are counted correctly \emph{always}; columns (3)--(4) use an indicator for saying politicians can identify how people voted \emph{never}. `BVR' equals one when the respondent's municipality had entered the biometric regime by 2021. `Bolsonaro approval' equals one when the respondent evaluates the executive as `good' or `very good'. All specifications control for low education, female, white, age, urban residence, and log municipality population. Standard errors clustered by municipality appear in parentheses. *** $p<0.01$, ** $p<0.05$, * $p<0.10$.
        \par}
    \end{minipage}
\end{table}

\textbf{Results.} \hyperref[tbl:lapop-backlash-2021]{Table~\ref*{tbl:lapop-backlash-2021}} shows two patterns. First, there is a clear Bolsonaro-alignment gap in perceptions of vote counting. In column (1), respondents who approve of Bolsonaro are about 22 percentage points less likely to say that votes are always counted correctly, even after conditioning on the available controls. In the interaction specification, the baseline approval gap is even larger in untreated municipalities: column (2) estimates a coefficient of $-0.324$ on `Bolsonaro approval' ($p<0.01$).

Second, the interaction does not support a backlash story in which BVR deepens that skepticism. For the vote-count outcome, the interaction term in column (2) is positive, about $0.198$, and marginally significant at conventional levels ($p=0.063$). Substantively, this means that the fair-count gap between Bolsonaro-aligned respondents and others is \emph{smaller}, not larger, in municipalities already exposed to BVR. Put differently, the cross-sectional evidence does not suggest that biometric exposure amplified distrust of the count among respondents aligned with Bolsonaro; if anything, it attenuated that distrust. For ballot secrecy, the estimates are much weaker. Columns (3)--(4) show little evidence that either BVR or the interaction between BVR and Bolsonaro approval systematically shifts whether respondents believe politicians can trace individual votes.

\textbf{Discussion.} These estimates should be interpreted cautiously. The analysis relies on a single public cross-section from 2021, so it is not a direct post-2022 election test and cannot distinguish pre-existing local differences from attitudinal responses to BVR itself. The individual alignment measure is also imperfect: approval of Bolsonaro's government is a reasonable proxy for partisan affinity, but it is not the same as verified vote choice. In addition, BVR is still measured at the municipality level rather than the respondent level, which raises the usual exposure-mismeasurement concerns. Finally, the integrity items were asked in a narrower module, which limits the available controls and keeps the estimation sample modest. Even with those caveats, the results are informative. They show a sizable Bolsonaro-alignment gap in trust that votes are counted correctly, but they do not show that BVR exposure intensified that gap. The most defensible reading is therefore a restrained one: the public-wave evidence points to skepticism among Bolsonaro-aligned respondents, but not to a clear BVR-specific backlash channel.
